\documentclass[10pt,a4paper]{amsart}
\usepackage[margin=19mm]{geometry}
\usepackage{amsmath,amssymb,mathtools}
\usepackage[scr=rsfs,cal=euler]{mathalfa}
\usepackage{xfrac}
\usepackage[unicode,hidelinks,bookmarks=false]{hyperref}
\newcommand{\ie}{i.e.\ }
\newcommand{\cf}{cf.\ }
\newcommand{\resp}{respectively\ }
\newcommand{\Subsection}{\subsection}
\providecommand{\nobreakdash}{\nobreak\hbox{-}}
\setlength{\emergencystretch}{2em}
\multlinegap0pt
\usepackage{bm}
\usepackage{bbm}
\usepackage{dsfont}
\usepackage{xspace}
\usepackage{cancel}
\usepackage{mathtools}
\usepackage{amsthm}
\makeatletter
\@ifundefined{lemma}{\newtheorem{lemma}{Lemma}}{}
\makeatother
\newcommand{\E}     {\mathbb{E}} 
\newcommand{\gk}[1]{\left\{#1\right\}}
\newcommand{\rk}[1]{\left(#1\right)}
\title[Phase Transitions in a Particle Model for the Self-Adaptive ]{Phase Transitions in a Particle Model for the Self-Adaptive Response to Cancer Dynamics}
\author{Christian Kuehn, Quirin Vogel}
\date{Source: arXiv:2508.03577v1 (CC BY 4.0)}
\begin{document}
\maketitle

\noindent\emph{Source and licence.} Phase Transitions in a Particle Model for the Self-Adaptive Response to Cancer Dynamics. Version arXiv:2508.03577v1 (CC BY 4.0), distributed under the Creative Commons Attribution 4.0 International licence, as recorded in the arXiv licence metadata for this exact version and shown on the abstract page https://arxiv.org/abs/2508.03577v1. Full rights notice: licenses/arxiv-2508.03577-CC-BY-4.0.txt.\par\smallskip

\noindent\textit{Editorial scope.} Counted unit: the Lemma whose source label is \texttt{eq:313251} in the article (source0.tex lines 556--565, source label \texttt{eq:313251}), delivered together with the article's own complete original proof of it (source0.tex lines 566--624, one \texttt{proof} environment of 59 lines). No mathematics is written, repaired or re-derived: statement and proof are the article's own text line for line. Required context retained uncounted: none. Every definition, display and label of those lines is retained unchanged. Omitted from this package and not redistributed: 1--555, 625--863. Nothing inside a retained excerpt is omitted or altered: every retained body is delivered exactly as the article's lines are written, so no omission receipt of any kind accompanies this package. Citations: the retained excerpts carry no citation command at all, so no bibliography entry of the article is transcribed and no reference is redistributed. Reference adaptation: none was needed and none was made; every reference inside the delivered unit resolves inside it, so no unresolved reference or citation is produced. Printed slips kept literal: no printed word, symbol, hypothesis or number of the counted statement or of its proof was altered. Layout and class: the article's own class is \texttt{article} with packages including graphicx, todonotes, pgfplots, subfigure, amsmath, amssymb, mathtools, graphics, color, verbatim, bbold, ulem, eucal, upgreek; the wrapper is the lane's pinned \texttt{amsart} editorial wrapper at 10pt on A4, which restates the theorem-like environments the retained text needs and adds the article's own macro definitions that the retained text needs (\texttt{\textbackslash E}, \texttt{\textbackslash gk}, \texttt{\textbackslash rk}), with extra wrapper packages (bm, bbm, dsfont, xspace, cancel, mathtools). The delivered numbering is the unit's own. Assets: no retained span contains a figure environment, a graphics-inclusion command, a PDF-page inclusion, a table or a TikZ picture; no image file is copied into this package, the package declares no asset dependency and no image file is redistributed with it. Licence: version arXiv:2508.03577v1 of this article is distributed under the Creative Commons Attribution 4.0 International licence, as recorded in the arXiv licence metadata for this exact version and shown on the abstract page https://arxiv.org/abs/2508.03577v1. A full rights notice is delivered at licenses/arxiv-2508.03577-CC-BY-4.0.txt. Characters: every retained excerpt is pure ASCII, so the delivered engine, XeLaTeX with its default Latin Modern text font, reports no missing character.
\begin{lemma}
    We have that
    \begin{equation}
        f(M-i)=\frac{M}{i!i\Gamma(a+1)q}\sum_{k=0}^{i-1}\frac{\Gamma(a+i-k)}{\Gamma(i-k+1)}\, .
    \end{equation}
    In particular
    \begin{equation}\label{eq:313251}
       f(0)=\E_0[H_M]=\frac{1}{q}\frac{\Gamma\rk{M+1+a}}{\Gamma(a+1)\Gamma(M+1)}\, .
    \end{equation}
\end{lemma}

\begin{proof}
Note that the following recursion holds for $f(i)$: using the Markov property of the chain and the linearity of the expectation, we get for $\in \gk{0,\ldots, M-1}$
    \begin{equation}\label{eq:recc}
\begin{split}
    f(M)&=0\, ,\\
    f(i)&=1+pf(0)+q\tfrac{i}{M}f(i)+q\rk{1-\tfrac{i}{M}}f(i+1)\, .
\end{split}
\end{equation}
If we make the ansatz
\begin{equation}\label{eq:Ansatz}
    1+pf(0)=b_{k-1}+pf(k)\, ,
\end{equation}
we can quickly derive the following facts:
\begin{enumerate}
    \item The recursion for $i=1$ can be rewritten as 
    \begin{equation}
        f(0)=\frac{1}{q}+f(1)\, .
    \end{equation}
    Multiplying with $p$ and adding $1$ yields $b_0=\frac{1}{q}$.
    \item If we plug Eq.~\eqref{eq:Ansatz} into Eq.~\eqref{eq:recc}, we obtain
    \begin{equation}
        f(k)=b_{k-1}+pf(k)+q\tfrac{k}{M}f(k)+q\rk{1-\tfrac{k}{M}}f(k+1)\, ,
    \end{equation}
    which is equivalent to
    \begin{equation}
        f(k)=\frac{b_{k-1}}{q\rk{1-\tfrac{k}{M}}}+f(k+1)\, .
    \end{equation}
    Multiplying with $p$ and adding $b_{k-1}$ yields the recursion
    \begin{equation}\label{eq:recbk}
        b_{k}=b_{k-1}+b_{k-1}\frac{p}{q\rk{1-\tfrac{k}{M}}}\, .
    \end{equation}
    \item The recursion ends with $1+pf(0)=b_{M-1}$\, .
\end{enumerate}
Eq.~\eqref{eq:recbk} can be rewritten as $b_{k}=b_{k-1}\rk{1+\tfrac{pM}{q(M-k)}}$. It is a linear multiplicative recursion in $(b_k)_k$ and hence solvable. We hence obtain that (with $a=pM/q$)
\begin{equation}\label{eq:valofhelp}
    1+pf(0)=\prod_{k=1}^M\rk{1+\frac{p}{q}\frac{M}{k}}=\frac{\prod_{k=1}^M\rk{k+\tfrac{Mp}{q}}}{M!}=\frac{\Gamma\rk{M+1+a}}{\Gamma(a+1)\Gamma(M+1)}\, .
\end{equation}
This immediately implies
\begin{equation}\label{eq:f0zerasympt}
    f(0)=\frac{\Gamma\rk{M+1+a}}{p\Gamma(a+1)\Gamma(M+1)}-\frac{1}{p}\sim \frac{(M+1)^aM}{a\Gamma(a+1)q}\sim \frac{M^{a+1}}{\Gamma(a+1)a}\, .
\end{equation}
Now, Eq.~\eqref{eq:recc} gives that
\begin{equation}
    f(M-i)=M\frac{1+pf(0)}{pM+qi}+\frac{qi}{pM+qi}f(M-i+1)\, .
\end{equation}
Inserting our result from Eq.~\eqref{eq:valofhelp} into the recursion gives that
\begin{equation}
    f(M-i)=M(1+pf(0))i!\sum_{k=0}^{i-1}\frac{q^{k}}{(i-k)!\prod_{j=0}^{k}\rk{pM+(i-j)q}}\, .
\end{equation}
Note that
\begin{equation}
    \prod_{j=0}^{k}\rk{pM+(i-j)q}=q^{k+1}\frac{\Gamma(a+i+1)}{\Gamma(a+i-k)}\, ,
\end{equation}
and hence
\begin{equation}
    f(M-i)=\frac{\Gamma(M+1+a)\Gamma(i+1)}{\Gamma(a+i+1)\Gamma(a+1)\Gamma(M)}\sum_{k=0}^{i-1}\frac{\Gamma(a+i-k)}{\Gamma(i-k+1)} =\frac{M}{i!i\Gamma(a+1)}\sum_{k=0}^{i-1}\frac{\Gamma(a+i-k)}{\Gamma(i-k+1)}\, .
\end{equation}
This concludes the proof.
\end{proof}

\end{document}
